\chapter{Limitations}

The methodology chapter already outlined the \emph{conceptual} limitations of the approach: the static environment and the cost-sensitive contextual bandit formulation, the binary label mapping, the fragility of the benchmark, and the scope of a single algorithm. This chapter complements them in light of the measured results, anchoring each limitation in concrete evidence rather than a general statement.

\section{Single Model and Seed}
The results correspond to a single model trained with a single seed. The reported bootstrap confidence intervals ($\pm 0.0002$ to $\pm 0.001$ at 95\%) quantify the sampling precision of the fixed test set, not the variance between training runs: they characterize how stable the metric is over resamplings of the same test set for \emph{this} model, but not how much performance would vary when retraining with different seeds. Estimating between-seed variance would require multiple full training runs, which have not been performed. The figures should therefore be read as the behavior of a well-characterized configuration point, not as an average over training randomness.

\section{Leakage in the Random Split}
The duplicate analysis shows that 22.30\% of the rows in the full dataset are exact duplicates and that 24.86\% of the test rows in the random split duplicate a training row (40.12\% among test attacks). Consequently, in-distribution metrics overestimate generalization capacity and should not be presented as the primary result in an operational sense. This limitation is methodologically mitigated by reporting Check C (day split) as the credible generalization signal, but it is not eliminated: the random figure remains the weakest rung of the project and is kept solely as evidence of learning capacity.

\section{Baseline Comparison}
The comparison with Random Forest was executed under the same protocol (canonical representation, scaling, balanced class weighting, and equivalent splits), with committed configuration and metric artifacts. Nevertheless, two limitations persist. First, the comparison covers a single supervised model; other competitive tabular classifiers, such as gradient-boosted ensembles, are discussed in the literature review but are not part of the evaluated protocol. Second, the baseline leave-one-CSV-out sweep (Wednesday reserved, attack recall 0.00719) does not yet have an equivalent committed QRDQN artifact, so on that rung the comparison is only of the baseline with itself and not between models. The comparison between QRDQN and Random Forest is direct and uses the same protocol on the fixed random split and the day split, but not on the leave-one-CSV-out split.

\section{Phase 2 External Validity}
The high Phase 2 performance (0.991862 accuracy) is obtained on traffic captured by the operator in a closed home lab and non-adversarial environment, from a single setting and with the generator's intent labels. It is not production traffic, it contains no real adversary, and it does not come from an independent source. Its external validity is therefore limited: a strong Phase 2 result is a controlled check that the feature contract and robust preprocessing transfer to a near-distribution capture, not a proof of real-world detection \parencite{SommerPaxson2010ClosedWorld,Layeghy2023CrossDomainNIDS}.

\section{Planned but Unexecuted Experiments}
Out of methodological honesty, the planned analyses whose artifacts have not been committed are explicitly listed. The learning curve protocol (training scale at different budgets) is planned but not executed: there is no learning curve artifact. The leave-one-CSV-out validation is implemented in code for QRDQN, but an aggregate metrics artifact has not been committed. Evaluation with multiple seeds to estimate between-run variance has not been carried out. None of these three capabilities are reported with figures in this thesis, and their absence is declared here rather than implied as completed.

\section{Future Work}
The above limitations define a coherent future work agenda: executing and committing the QRDQN learning curve and leave-one-CSV-out; retraining with multiple seeds to quantify variance; extending the comparison to other supervised models; and, above all, validating on independent labeled traffic from multiple environments before considering any step towards active deployment, which in turn would require a specific security analysis. These extensions strengthen the evidence without altering the experimental architecture already established.
